\documentclass{article}
\usepackage[paperwidth=841mm,paperheight=1189mm,margin=0mm]{geometry}
\usepackage{tcolorbox}
\tcbuselibrary{poster,skins}
\usepackage{graphicx}
\usepackage{amsmath,amssymb}
\usepackage{enumitem}
\usepackage[table]{xcolor}
\usepackage{lmodern}
\usepackage[T1]{fontenc}
\usepackage{array}
\pagestyle{empty}

\definecolor{primary}{HTML}{12231C}
\definecolor{secondary}{HTML}{C67B3B}
\definecolor{accent}{HTML}{31A688}
\definecolor{amber}{HTML}{E3B24E}
\definecolor{bgposter}{HTML}{EDE3D0}
\definecolor{greencard}{HTML}{EEF7EF}
\definecolor{stonecard}{HTML}{F3F1EA}
\definecolor{sandcard}{HTML}{FFF2DA}
\definecolor{darkcardbg}{HTML}{E3EDDF}
\definecolor{greentitlebg}{HTML}{DBEEE1}
\definecolor{stonetitlebg}{HTML}{E7E3D8}
\definecolor{sandtitlebg}{HTML}{F5DFB5}
\definecolor{textdark}{HTML}{101820}
\definecolor{textmuted}{HTML}{4B5A52}

\pagecolor{bgposter}
\color{textdark}

\setlength{\parindent}{0pt}
\setlist[itemize]{leftmargin=26pt, itemsep=7pt, parsep=1pt, topsep=2pt,
  label={\color{accent}$\blacktriangleright$}}

\newcommand{\bodyfont}{\fontsize{31}{39}\selectfont}
\newcommand{\smallbody}{\fontsize{26}{33}\selectfont}
\newcommand{\captionfont}{\fontsize{23}{29}\selectfont}
\newcommand{\statnumber}[1]{{\fontsize{62}{74}\selectfont\bfseries\color{primary}#1}}
\newcommand{\statlabel}[1]{{\fontsize{24}{30}\selectfont\color{textdark}#1}}
\newcommand{\posterfig}[3]{%
  \centering\includegraphics[width=#1\linewidth]{#2}\\[2mm]
  {\captionfont\color{textmuted}\textit{#3}}\vspace{1mm}%
}
\newcommand{\callout}[2]{%
  \colorbox{#1!22}{\parbox{0.955\linewidth}{\smallbody #2}}%
}

\tcbset{
  basecard/.style={
    enhanced, arc=0pt, boxrule=0pt,
    left=18pt, right=18pt, top=6pt, bottom=6pt,
    fonttitle=\fontsize{37}{45}\selectfont\bfseries,
    toptitle=6pt, bottomtitle=6pt,
    titlerule=2pt,
    valign=top,
    drop shadow={opacity=0.16},
  },
  greencard/.style={
    basecard, colback=greencard,
    borderline west={6pt}{0pt}{accent},
    coltitle=primary, colbacktitle=greentitlebg,
    titlerule style={accent!55},
  },
  stonecard/.style={
    basecard, colback=stonecard,
    borderline west={6pt}{0pt}{primary},
    coltitle=primary, colbacktitle=stonetitlebg,
    titlerule style={primary!45},
  },
  sandcard/.style={
    basecard, colback=sandcard,
    borderline west={6pt}{0pt}{secondary},
    coltitle=secondary, colbacktitle=sandtitlebg,
    titlerule style={secondary!55},
  },
  highlightcard/.style={
    basecard, colback=darkcardbg,
    borderline west={7pt}{0pt}{primary},
    coltitle=white, colbacktitle=primary,
  }
}

\begin{document}
\begin{tcbposter}[
  coverage={spread},
  poster={columns=2, rows=24, spacing=0mm},
]

\posterbox[
  enhanced, colback=primary, colframe=primary, colupper=white,
  arc=0pt, boxrule=0pt,
  left=36pt, right=36pt, top=10pt, bottom=8pt,
  halign=center, valign=center,
  drop shadow={opacity=0.24}
]{name=title, column=1, span=2, between=top and row3}{
  {\fontsize{70}{84}\selectfont\bfseries SpriteCraft: Texture Style Transfer for Minecraft Resource Packs}\\[10pt]
  {\fontsize{33}{41}\selectfont Diffusion generation + structure-preserving recolor fallback for 32$\times$32 block textures}\\[6pt]
  {\fontsize{27}{34}\selectfont Ray Zhang \quad Shanghai Jiao Tong University}
}

\posterbox[
  enhanced, colback=primary!12!white, boxrule=0pt, arc=0pt,
  left=10pt, right=10pt, top=6pt, bottom=6pt,
  valign=center, borderline south={3pt}{0pt}{primary!35}
]{name=stats, column=1, span=2, between=row3 and row5}{
  \centering
  \begin{minipage}[c]{0.235\linewidth}\centering
    \statnumber{32$\times$32}\\[-1pt]\statlabel{canonical RGB texture grid}
  \end{minipage}\hfill
  \begin{minipage}[c]{0.235\linewidth}\centering
    \statnumber{7.7M+1.0M}\\[-1pt]\statlabel{diffusion + recolor parameters}
  \end{minipage}\hfill
  \begin{minipage}[c]{0.235\linewidth}\centering
    \statnumber{40}\\[-1pt]\statlabel{cosine DDPM sampling steps}
  \end{minipage}\hfill
  \begin{minipage}[c]{0.235\linewidth}\centering
    \statnumber{4}\\[-1pt]\statlabel{\texttt{run33} fully trained packs}
  \end{minipage}
}

\posterbox[greencard, title={Problem: Small Tiles Amplify Every Error}
]{name=problem, column=1, between=row5 and row9}{
  \bodyfont
  \begin{itemize}
    \item Resource packs require hundreds of hand-made blocks.
    \item Low resolution magnifies hue and edge mistakes.
    \item Vanilla filenames give paired pack supervision.
  \end{itemize}
  \vspace{2mm}
  \callout{accent}{\textbf{Task:} preserve material identity, adopt pack palette, and remain plausible as a repeating Minecraft tile.}
}

\posterbox[stonecard, title={Pipeline: Vanilla-Anchored Per-Pack Training}
]{name=pipeline, column=1, between=row9 and row14}{
  \posterfig{0.965}{figures/spritecraft-pipeline-overview.png}{Raw packs are filtered to valid full-cube block textures, normalized to 32$\times$32, paired by filename, and exported as per-pack datasets.}
}

\posterbox[sandcard, title={Method: Expressive Diffusion + Safe Recolor}
]{name=architecture, column=1, between=row14 and row20}{
  \posterfig{0.965}{figures/spritecraft-dual-model-architecture.png}{StyleAwareUNet samples expressive candidates; RecolorNet supplies structure-preserving residual color transfer; the selector chooses the final output.}
}

\posterbox[highlightcard, title={Training and Routing Rules}
]{name=methodrules, column=1, between=row20 and bottom}{
  \bodyfont
  \begin{itemize}
    \item SNR-weighted $\hat{x}_0$ reconstruction and content losses.
    \item Source-relative structure, contrast, hue, color moments.
    \item Candidate score = structural quality + style quality.
    \item Recolor fallback for low-quality diffusion outputs.
  \end{itemize}
}

\posterbox[sandcard, title={Pack-Level Results at \texttt{run33} Step 10k}
]{name=results, column=2, between=row5 and row10}{
  \smallbody
  \renewcommand{\arraystretch}{1.35}
  \setlength{\tabcolsep}{7pt}
  \centering
  \begin{tabular}{>{\raggedright\arraybackslash}p{0.28\linewidth}ccc}
    \rowcolor{secondary!20}
    Pack & Diff. & Recolor & Selected \\
    Ashen 16x & .107 / .315 & .110 / .383 & \textbf{.110 / .389} \\
    Bare Bones & \textbf{.059 / .560} & .068 / .559 & .062 / .526 \\
    Cartoon & \textbf{.136 / .309} & .160 / .264 & .151 / .225 \\
    Chibli 64x & \textbf{.126 / .302} & .188 / .189 & .188 / .189 \\
  \end{tabular}\\[3mm]
  {\captionfont\color{textmuted}Cells show MAE / pixel accuracy. The selector is conservative: it optimizes plausible tile structure, not target-aware MAE alone.}
}

\posterbox[greencard, title={Qualitative Result: One Texture, Four Pack Styles}
]{name=oak, column=2, between=row10 and row15}{
  \posterfig{0.93}{figures/run33-multi-pack-oak-planks.png}{Selected \texttt{oak\_planks.png} outputs preserve plank layout while moving palette and contrast toward each pack.}
}

\posterbox[stonecard, title={Routing Stress Test: Fine Structure Wins}
]{name=routing, column=2, between=row15 and row20}{
  \posterfig{0.93}{figures/run33-ashen-routing.png}{Bookshelf and TNT expose diffusion's structural risk. Recolor fallback keeps symbols and narrow repeated structure more legible.}
}

\posterbox[highlightcard, title={Takeaways and Next Steps}
]{name=takeaways, column=2, between=row20 and bottom}{
  \bodyfont
  \begin{itemize}
    \item Dual models are safer than diffusion alone.
    \item Easy packs accept more diffusion detail.
    \item Hard packs need conservative structure gates.
    \item Future: alpha channels, learned support ranking, multi-resolution refinement, learned routing thresholds.
  \end{itemize}
  \vspace{2mm}
  \callout{amber}{\textbf{Best current fit:} vanilla-adjacent, minimalist, and medium-complexity stylized Minecraft resource packs.}
}

\end{tcbposter}
\end{document}

